% lerobot-bench --- 4-page arXiv writeup (cs.RO primary, cs.LG secondary).

\documentclass[10pt]{article}

\usepackage[a4paper,margin=0.75in]{geometry}
\usepackage[utf8]{inputenc}
\usepackage[T1]{fontenc}
\usepackage{lmodern}
\usepackage{microtype}
\usepackage{amsmath,amssymb}
\usepackage{booktabs}
\usepackage{subcaption}
\usepackage{graphicx}
\usepackage{xcolor}
\usepackage{hyperref}
\usepackage{url}
\usepackage[compact]{titlesec}
\titlespacing*{\section}{0pt}{1.4ex plus 0.3ex minus 0.2ex}{0.8ex plus 0.2ex}
\titlespacing*{\subsection}{0pt}{1.0ex plus 0.2ex minus 0.2ex}{0.5ex plus 0.1ex}

\hypersetup{
    colorlinks=true,
    linkcolor=blue!50!black,
    citecolor=blue!50!black,
    urlcolor=blue!50!black,
}

\title{lerobot-bench: a multi-policy reproducible evaluation\\
suite for pretrained LeRobot policies}
\author{Th\'eo Hermann\\
\texttt{thermann.ai@gmail.com} \\
\texttt{https://github.com/thrmnn/lerobot-bench}}
\date{\today}

\begin{document}

\maketitle

% ---------------------------------------------------------------- %
% Abstract                                                         %
% ---------------------------------------------------------------- %

\begin{abstract}
Reported success rates for pretrained LeRobot~\cite{cadene2024lerobot}
policies live in self-reported tables across heterogeneous protocols,
with no shared seed accounting or confidence intervals. We introduce
\emph{lerobot-bench}, a public reproducible evaluation suite that runs
a fixed roster of six policies --- four learned (Diffusion Policy, ACT,
SmolVLA, XVLA) and two weights-free floor baselines (no-op,
uniform-random) --- on six LeRobot sim envs (PushT, Aloha-transfer-cube,
and the four LIBERO suites~\cite{liu2023libero}), 22~runnable
\textit{(policy, env)} cells after the checkpoint-availability filter.
Each cell runs a multi-seed contract (5~seeds $\times$ 50~episodes,
$N{=}250$ binary outcomes) with Wilson and bootstrap 95\% confidence
intervals plus paired comparisons under explicit
minimum-detectable-effect (MDE) bounds, and every entry is pinned to a
fixed revision SHA so a parquet row is replayable bit-for-bit. The
headline finding is a bug we caught in our \emph{own} harness: an early
build silently skipped dataset normalization on \texttt{observation.images.top},
feeding ACT un-normalized images and pinning ACT on Aloha-transfer-cube at
0.016~[0.006, 0.040] --- below the env's random floor. After we fixed the
normalization (PR~\#51) the same checkpoint reads 0.824~[0.772, 0.866]; a
controlled $2{\times}2$ ablation (normalization $\times$ inference settings)
shows the recovery is entirely the norm fix, while temporal ensembling is a
wash. The lesson is that the harness itself, not the architecture, was the
load-bearing variable, and that an eval suite must be audited as carefully as
the policies it scores. SmolVLA on LIBERO-10 measures
$\hat{p}{=}0.252$~[0.202, 0.309] under our v1 default protocol vs.\ the
0.71 reported~\cite{shukor2025smolvla}, but a v1.0.1 audit bounds the
gap by scope mismatch: the paper averages 10~tasks $\times$ 10~trials,
we ran \texttt{task\_id=0} $\times$ 5~seeds $\times$ 50~episodes, and
74.8\% of failed \texttt{libero\_10} episodes hit the 520-step cap vs.\
canonical LIBERO's 600. These numbers are real for the scope measured;
the audit constrains what they mean. The implication is that single-seed
self-reported numbers should not be compared head-to-head without
re-running under a shared protocol.
\end{abstract}

% ---------------------------------------------------------------- %
% Introduction                                                     %
% ---------------------------------------------------------------- %

\section{Introduction}

LeRobot~\cite{cadene2024lerobot} ships pretrained checkpoints for
Diffusion Policy~\cite{chi2023diffusion}, ACT~\cite{zhao2023learning},
SmolVLA~\cite{shukor2025smolvla}, XVLA~\cite{bu2025xvla}, and the
$\pi_0$ family~\cite{black2024pi0} across simulated manipulation envs
(PushT, Aloha, LIBERO). Individual policy papers report success rates
on these envs, but the protocols differ in seed counts, episode
budgets, success thresholds, and how confidence intervals (if any) are
derived. The result is a set of numbers in PDFs that cannot be trusted
head-to-head --- a problem documented repeatedly in adjacent
fields~\cite{henderson2018deep,agarwal2021deep}.

We fill that gap with \emph{lerobot-bench}, released across three
surfaces: open data (a public Hub dataset
\texttt{thrmnn/lerobot-bench-v1} of every per-episode outcome plus an
MP4 of every rollout, indexed by \textit{(policy, env, seed, episode)});
open eval (a pinned, scriptable pipeline that re-derives any parquet row
in $<1$\,h on commodity hardware, plus the resumable orchestrator that
ran the full sweep); and open analysis (the notebook
\texttt{notebooks/01-write-finding.ipynb} that produces every table and
figure here). The v1 roster is four learned policies and two
weights-free baselines crossed with six envs; after filtering to cells
with a public Hub checkpoint the matrix has 22~runnable cells. The
$\pi_0$ family is deferred to v1.1 for a host-RAM reason detailed in the
Limitations.

The contribution is methodological as much as empirical: a seeding
contract that makes per-cell numbers bit-reproducible at cell
boundaries, a bootstrap protocol that makes CIs defensible at $N{=}250$
binary outcomes per cell, and an MDE bound that prevents over-claiming
on small deltas. The benchmark itself is the data; this paper is the
defense.

% ---------------------------------------------------------------- %
% Related Work                                                     %
% ---------------------------------------------------------------- %

\section{Related work}
\label{sec:related}

\paragraph{Robot-learning benchmarks.}
PushT~\cite{chi2023diffusion} and the Aloha simulation
suite~\cite{zhao2023learning} were each introduced as harnesses for a
single policy class, and the success rates in those papers remain the de
facto reference. LIBERO~\cite{liu2023libero} was designed as a benchmark
from the outset, with four suites (\textit{spatial}, \textit{object},
\textit{goal}, \textit{long}; the latter exposed in lerobot as
\texttt{libero\_10}) targeting compositional generalisation.
RLBench~\cite{james2020rlbench} and CALVIN~\cite{mees2022calvin} offer
larger task taxonomies but host no pretrained LeRobot-compatible
checkpoints; Open X-Embodiment~\cite{bohg2025open} provides
cross-embodiment data but no evaluation protocol on these envs. None
ship a multi-policy leaderboard with shared seed accounting, multi-seed
CIs, and re-runnable rows for the LeRobot stack --- the gap we fill.

\paragraph{Statistical rigor in policy evaluation.}
Henderson et al.~\cite{henderson2018deep} showed single-seed deep RL
claims are largely irreproducible across seeds; Agarwal et
al.~\cite{agarwal2021deep} formalised multi-seed reporting with
stratified bootstrap CIs and per-task distributional summaries. We adopt
the same posture --- multi-seed evaluation, bootstrap CIs, paired
comparisons, and MDE bounds --- for the LeRobot manipulation matrix,
where no comparable protocol has been published.

\paragraph{Floor baselines.}
A leaderboard without no-op and uniform-random baselines cannot separate
``the env is hard'' from ``the policy is good'': some LeRobot envs admit
non-zero success under random action. We therefore evaluate two
weights-free baselines on every cell as a leaderboard floor, mirroring
the OpenAI Gym~\cite{brockman2016openai} convention from which the
LeRobot env API descends.

% ---------------------------------------------------------------- %
% Methods                                                          %
% ---------------------------------------------------------------- %

\section{Methods}
\label{sec:methods}

\paragraph{Sweep contract.}
The unit of evaluation is a \textit{cell}:
$(\text{policy}, \text{env}, \text{seed\_idx})$ with
$\text{seed\_idx} \in \{0,\ldots,4\}$. Each cell runs $n_\text{ep}$
episodes; the contract is $n_\text{ep}{=}50$ unless the auto-downscope
rule (below) reduces it, giving $N = 5 \cdot 50 = 250$ binary outcomes.
An episode runs from \texttt{env.reset(seed=$s$)} until
\texttt{terminated}, \texttt{truncated}, or step count reaches
env-specific \texttt{max\_steps} (PushT 300, Aloha 400, LIBERO
$\{280,280,300,520\}$ for spatial/object/goal/10, inherited from
lerobot defaults; canonical LIBERO uses 600 for all four, an audit
caveat in Section~\ref{sec:results-audit}). Success is the per-env
\texttt{SUCCESS\_REWARD} threshold; the episode succeeds iff its
final-step reward meets or exceeds it. PR~\#90 adds a selectable
\texttt{--canonical} criterion (paper termination rule, 600-step LIBERO
cap); v1 ships the lerobot-default rule and v1.1 reruns affected cells.

\paragraph{Seeding contract.}
Per cell, three RNGs are seeded with
$\text{seed}_\text{cell} = \text{seed\_idx} \cdot 1000$
(\texttt{numpy}, \texttt{torch}, \texttt{torch.cuda}). Per episode the
env is reset with $\text{seed}_\text{ep}(e) = \text{seed}_\text{cell}+e$.
Policy stochasticity inherits the torch generator and is \emph{not}
re-seeded per episode, so the generator advances across episodes within
a cell: mid-cell resume is \emph{not} bit-reproducible and a crashed
cell restarts from $e{=}0$. The orchestrator
(\texttt{checkpointing.plan\_resume}) classifies cells as completed,
partial, or pending and re-runs partial cells from scratch.

\paragraph{Per-cell confidence intervals.}
For each cell we report $\hat{p} = (\text{successes})/N$ with two CIs we
sanity-check against each other: (i)~the closed-form Wilson score
interval~\cite{wilson1927probable}, recommended for proportions near 0
or 1~\cite{agresti1998approximate}; and (ii)~a percentile-bootstrap 95\%
interval~\cite{efron1979bootstrap,efron1993introduction} with
$n_\text{resamples}{=}10{,}000$ and a seeded
\texttt{numpy.random.Generator(42)}. The resampling unit is the
\emph{episode}, not the seed: bootstrapping over 5 seeds gives wide CIs
and conflates between-seed variance with the inferential question. The
two intervals agree to within 1\,pp at $N{=}250$ across the $\hat{p}$
grid; we show Wilson in the table (closed-form, no RNG) and bootstrap in
the forest plot.

\paragraph{Paired comparisons.}
For cell pairs with identical $(\text{seed\_idx}, e)$ pairings (same
protocol, same env) we report a paired $\Delta\hat{p}$ two ways: a
paired bootstrap (\texttt{stats.paired\_delta\_bootstrap}, same
$n_\text{resamples}$ and seed) giving a pivotal CI on the
difference~\cite{efron1993introduction}, and the two-sided Wilcoxon
signed-rank test~\cite{wilcoxon1945individual} on per-episode
differences (\texttt{zero\_method=`wilcox'}, which for binary outcomes
reduces in the limit to McNemar's test~\cite{mcnemar1947note}). Effect
size is Cohen's $h$~\cite{cohen1988statistical},
$h = 2\arcsin(\sqrt{p_1}) - 2\arcsin(\sqrt{p_2})$, with bands
$|h|<0.2$ small, $<0.5$ medium, $\geq 0.8$ large. When an env carries
more than one ranking claim, the family-wise rate is controlled by
Holm--Bonferroni~\cite{bonferroni1936teoria} at $\alpha{=}0.05$ over the
$k$ within-env comparisons.

\paragraph{Minimum detectable effect.}
At $N{=}250$ the Wilson 95\% half-width at worst case $\hat{p}{=}0.5$ is
$\text{HW}=0.0615$; any $|\Delta\hat{p}| < 2\cdot\text{HW} = 0.1230$ is
inconclusive at this $N$ regardless of $p$-value, since two cells with
half-widths up to 0.0615 can produce a smaller sample delta from
sampling noise alone. The threshold shrinks for saturated/near-zero
cells ($0.0543$ at $\hat{p}{=}0.95$); \texttt{docs/MDE\_TABLE.md}
tabulates the per-$\hat{p}$ threshold and the empirical paired bootstrap
MDE at $\rho\in\{0,0.3,0.5\}$. Every paired comparison in
Section~\ref{sec:results} is annotated against this threshold; deltas
below it are reported as \emph{inconclusive at $N{=}250$}, not as
findings.

\paragraph{Auto-downscope rule.}
\label{sec:downscope}
The per-cell wall-clock budget is 3~hours on a single RTX~4060 laptop
GPU (8\,GB VRAM). Given a per-policy mean step latency $\overline{t}$
(ms/step) from Day-0 calibration, the episodes-per-seed count is
\begin{equation}
    n_\text{ep} = \min\!\left(50,\ \left\lfloor
        \frac{3 \cdot 3600}{5 \cdot \texttt{env.max\_steps} \cdot
        \overline{t} / 1000} \right\rfloor\right).
\end{equation}
Seed count is fixed at five; only $n_\text{ep}$ shrinks, so a
down-scoped cell lands at $N\in\{50,100\}$ with a wider CI flagged in
the table caption. Two v1 cells were down-scoped. If the per-cell
minimum ($N{=}50$) still exceeds budget, the policy is dropped from v1,
not silently truncated.

\paragraph{Sparse matrix and reproducibility.}
A \textit{(policy, env)} cell is N/A only when no public Hub checkpoint
exists at lock-in, reducing the $6\times6$ grid to 22~runnable cells;
sparsity is a visible gap in Table~\ref{tab:leaderboard}, not silently
skipped, and cross-policy comparisons are made only within shared envs.
Each parquet row carries its code commit SHA, the lerobot pin
(\texttt{lerobot==0.5.1}), and the SHA-256 of its MP4. Every entry in
\texttt{configs/policies.yaml} is locked to a fixed
\texttt{revision\_sha} (resolved from
\texttt{HfApi().model\_info(repo\_id).sha} at lock-in, 2026-05-03); the
eval loop refuses to load a policy without one. A single row replays via
\texttt{python scripts/run\_one.py --policy P --env E --seed S}. The
full sweep took $\approx 25$~hours of GPU time on the reference hardware
(1$\times$ RTX 4060 Laptop, WSL2, lerobot 0.5.1; wall-clock
2026-05-22T00:35Z $\to$ 2026-05-23T01:39Z).

% ---------------------------------------------------------------- %
% Results                                                          %
% ---------------------------------------------------------------- %

\section{Results}
\label{sec:results}

\paragraph{Per-policy leaderboard.}
Table~\ref{tab:leaderboard} reports pooled per-policy success rate
$\hat{p}$ across supported envs with the Wilson 95\% interval; rows with
$N{<}250$ are auto-downscoped (Section~\ref{sec:downscope}) and flagged
$\dagger$. The full per-\textit{(policy, env)} breakdown is
Table~\ref{tab:breakdown}.

\begin{table}[t]
\centering
\small
\begin{minipage}[t]{0.40\linewidth}
\centering
\begin{tabular}{lrrl}
\toprule
Policy & $N$ & $\hat{p}$ & 95\% CI \\
\midrule
ACT              & 250  & 0.824 & [0.772, 0.866] \\
Diffusion        & 125$\dagger$ & 0.816 & [0.739, 0.874] \\
SmolVLA          & 1000 & 0.621 & [0.591, 0.651] \\
random           & 1500 & 0.009 & [0.005, 0.015] \\
no-op            & 1500 & 0.000 & [0.000, 0.003] \\
\bottomrule
\end{tabular}
\subcaption{Pooled per-policy.}
\label{tab:leaderboard}
\end{minipage}\hfill
\begin{minipage}[t]{0.57\linewidth}
\centering
\begin{tabular}{llrrl}
\toprule
Policy & Env & $N$ & $\hat{p}$ & 95\% CI \\
\midrule
ACT       & aloha\_xfer    & 250 & 0.824 & [0.772, 0.866] \\
Diffusion & pusht          & 125$\dagger$ & 0.816 & [0.739, 0.874] \\
SmolVLA   & libero\_goal   & 250 & \textbf{0.928} & [0.889, 0.954] \\
SmolVLA   & libero\_spat.  & 250 & 0.776 & [0.720, 0.823] \\
SmolVLA   & libero\_obj.   & 250 & 0.528 & [0.466, 0.589] \\
SmolVLA   & libero\_10     & 250 & \textbf{0.252} & [0.202, 0.309] \\
random    & aloha\_xfer    & 250 & 0.052 & [0.031, 0.087] \\
random    & other envs     & 1250 & 0.000 & [0.000, 0.003] \\
no-op     & all envs       & 1500 & 0.000 & [0.000, 0.003] \\
\bottomrule
\end{tabular}
\subcaption{Per-\textit{(policy, env)}.}
\label{tab:breakdown}
\end{minipage}
\caption{Success rate with Wilson 95\% CI. (a) pooled per-policy across
supported envs; (b) per-cell breakdown, $N{=}5\cdot n_\text{ep}$ binary
outcomes per cell. $\dagger$: auto-downscoped
diffusion$\times$pusht cell ($n_\text{ep}{=}25$, $N{=}125$;
Section~\ref{sec:downscope}). \texttt{random} on
\texttt{aloha\_transfer\_cube} is the only above-floor baseline (the
\texttt{gym-aloha} heuristic is the most permissive in the matrix). MDE
band at $N{=}250$, $\hat{p}{=}0.5$ is $0.123$; 22~runnable cells; XVLA
deferred.}
\end{table}

\paragraph{Headline: the harness, not the architecture, was the
load-bearing variable --- a normalization bug we caught in our own eval.}
The single most consequential v1 result began as ACT's apparent failure
on \texttt{aloha\_transfer\_cube} ($\hat{p}{=}0.016$~[0.006, 0.040], below
the env's random floor) and ended as a bug in \emph{our} harness. An early
build of the eval loop silently failed to apply dataset normalization stats
to \texttt{observation.images.top}, feeding ACT un-normalized image
observations; PR~\#51 fixed it
(\texttt{\_recover\_dataset\_stats\_from\_safetensors} now disambiguates
buffer names against the config's \texttt{feature\_keys}). After the fix the
same checkpoint reads the canonical \textbf{0.824~[0.772, 0.866]}
(Table~\ref{tab:breakdown}). To attribute the recovery cleanly we ran a
controlled $2{\times}2$ ablation crossing normalization \{buggy, fixed\}
with inference \{Hub-default, paper-settings\}, each cell $N{=}250$
(Table~\ref{tab:act-ablation}). The recovery is 100\% the normalization fix:
on broken normalization, switching to the paper's inference settings does
nothing ($0.016 \to 0.016$), and on fixed normalization, Hub-default vs.\
paper settings are statistically indistinguishable (0.812 vs.\ 0.768,
overlapping Wilson CIs). \emph{Temporal ensembling is a wash, not the
cause.} The lesson generalises: an eval suite must be audited --- normalization,
buffer wiring, observation plumbing --- as carefully as the policies it
scores (Section~\ref{sec:results-audit}).

\begin{table}[t]
\centering
\small
\begin{tabular}{lll}
\toprule
& Hub-default inference & Paper-settings inference \\
\midrule
Normalization buggy & 0.016 & 0.016 \\
Normalization fixed & 0.812 & 0.768 \\
\bottomrule
\end{tabular}
\caption{ACT $\times$ \texttt{aloha\_transfer\_cube}, controlled
$2{\times}2$ ablation (normalization $\times$ inference settings), each
cell $N{=}250$ (5~seeds $\times$ 50~episodes). The entire recovery is the
normalization fix (PR~\#51): broken normalization stays at 0.016
regardless of inference settings, and once normalization is fixed the
Hub-default and paper inference settings are statistically
indistinguishable (overlapping Wilson CIs), so temporal ensembling is a
wash. The fixed+Hub-default cell (0.812) is the same condition as the
canonical leaderboard number 0.824~[0.772, 0.866] (Table~\ref{tab:breakdown})
measured in a separate $N{=}250$ run; the two agree within CI. Source:
\texttt{scripts/probes/probe\_act\_normalization\_ablation.py}.}
\label{tab:act-ablation}
\end{table}

\paragraph{SmolVLA on LIBERO under our protocol.}
SmolVLA's pooled rate across the four LIBERO suites is
0.621~[0.591, 0.651] under our protocol (\texttt{task\_id=0}, 5~seeds
$\times$ 50~episodes, env-default caps), with a sharp per-suite
gradient: 0.928~[0.889, 0.954] (\texttt{goal}) and 0.776~[0.720, 0.823]
(\texttt{spatial}) against 0.528~[0.466, 0.589] (\texttt{object}) and
0.252~[0.202, 0.309] (\texttt{libero\_10}). Published SmolVLA-0.45B
numbers~\cite{shukor2025smolvla} are suite averages over 10~tasks
$\times$ 10~trials (0.90/0.96/0.92/0.71 for the same suites): ours is a
single-task probe at \texttt{task\_id=0}, not an apples-to-apples
replication, and \texttt{libero\_10} is additionally a lower bound at
the 520-step cap (74.8\% of failed episodes truncated; canonical LIBERO
uses 600). The within-protocol measurement is tight (Wilson half-width
$0.054$ at $N{=}250$); the replication claim defers to v1.1
(Section~\ref{sec:results-audit}). On \texttt{libero\_object} the
across-seed variance is $1{\times}10^{-4}$ (five seeds agree on
$\hat{p}{=}0.528$ to within one episode) vs.\ $5.9{\times}10^{-3}$ on
\texttt{spatial} and $2.7{\times}10^{-3}$ on \texttt{libero\_10}, so this
failure is consistent across seeds --- a structural cap on a subset of
object configurations, not a sample-size problem.

\paragraph{Baseline sanity.}
The no-op baseline succeeds zero times across all six envs ($N{=}1500$,
Wilson upper bound $0.003$). Uniform-random succeeds zero times on
\texttt{pusht} and every LIBERO suite and is the only above-floor
baseline on \texttt{aloha\_transfer\_cube} ($\hat{p}{=}0.052$,
[0.031, 0.087]), where the \texttt{gym-aloha} success heuristic is the
most permissive in the matrix. The floor is therefore not a single
number across the leaderboard --- exactly why we run it on every cell.


\paragraph{Methodology caveats (v1.0.1 audit).}
\label{sec:results-audit}
We ran a static audit against each policy's source paper and each env's
canonical protocol; three confirmed mismatches constrain what the
headline cells mean without invalidating the measurements.
(1)~\emph{Task scope.} The SmolVLA paper~\cite{shukor2025smolvla}
Table~2 averages 10~tasks $\times$ 10~trials per suite; v1 runs
\texttt{task\_id=0} $\times$ 5~seeds $\times$ 50~episodes, so every
SmolVLA paper-vs-measured delta is a single-task envelope claim, not a
replication (PR~\#84). (2)~\emph{ACT normalization bug (self-caught,
fixed).} An early build of our eval loop silently skipped applying
dataset normalization stats to \texttt{observation.images.top}, feeding
ACT un-normalized images and pinning it at \textbf{0.016~[0.006, 0.040]}.
PR~\#51 fixed it. The controlled $2{\times}2$ ablation
(Table~\ref{tab:act-ablation}) attributes the entire recovery to the
normalization fix --- on broken normalization, paper inference settings
change nothing ($0.016 \to 0.016$); on fixed normalization, Hub-default
and paper settings are indistinguishable (0.812 vs.\ 0.768, overlapping
Wilson CIs) --- so temporal ensembling is a wash, not the cause. We had
earlier attributed the recovery to inference settings; the ablation
corrects that. (3)~\emph{LIBERO
step caps (probe resolved for \texttt{libero\_10}).} Canonical
LIBERO~\cite{liu2023libero} uses 600~steps for all four suites; v1
inherited lerobot's caps \{280,280,300,520\}, with failed-episode
cap-hit rates 22.4\%, 47.2\%, 7.2\%, 74.8\%. Re-running
\texttt{libero\_10} at cap=600 measures \textbf{0.256~[0.206, 0.314]}
vs.\ 0.252 ($\Delta{=}+0.4$\,pp, within either half-width); cap-hits
stay at 74.4\% at 600, so \texttt{libero\_10} is
\emph{policy-bottlenecked}, not cap-bottlenecked. The other three suites
remain provisional lower bounds pending cap=600 reruns in v1.1 (PR~\#90
ships a selectable \texttt{--canonical} criterion). No patch rewrites any
parquet row: v1 numbers stay valid for the scope they were measured
under, and the audit constrains the comparative claim, not the
measurement.

\paragraph{XVLA deferral.}
\texttt{xvla\_libero} ran across all four LIBERO suites (875~episodes)
but is \textbf{deferred from the v1 leaderboard}. Two Hub-artifact wiring
bugs (missing rotation-6D$\to$axis-angle post-processor and ImageNet
normalization) were patched (PRs~\#71/\#74), yet a third unresolved issue
still produced 0/40 in a post-patch sanity sweep (Wilson upper bound
$0.088$ at $N{=}40$, consistent with a non-functional load);
\texttt{docs/DEFERRED\_POLICIES.md} documents all three and XVLA returns
in v1.1.

\paragraph{Stats sanity and failure taxonomy.}
\label{sec:taxonomy}
Wilson and pivotal-bootstrap 95\% CIs agree to within one half-width on
every cell; no cell triggers the disagreement flag. The orchestrator
advanced through 13 distinct \texttt{code\_sha} values across the 110
\textit{(policy, env, seed)} cells, but the revision pins and
\texttt{lerobot\_bench.stats} helpers were constant (only orchestrator
and env-glue drifted) and every row carries its SHA for audit. We also
hand-label 5--10 failed rollouts per cell into six modes
(\texttt{docs/FAILURE\_TAXONOMY.md}: trajectory overshoot, gripper slip,
timeout, wrong object, premature release, drift), taking the first that
fits chronologically and stratifying one draw per seed; the
\texttt{libero\_object} cluster localises to wrong-object on a fixed
object subset, consistent with the structural cap above.

% ---------------------------------------------------------------- %
% Discussion                                                       %
% ---------------------------------------------------------------- %

\section{Discussion}

\paragraph{What the leaderboard says.}
Three things, in order of confidence. First, the protocol works: every
row is bit-replayable, Wilson and bootstrap CIs agree to within one
half-width on every cell, and 110~cells dispatched with zero failures
over 25~hours unattended. Second, harness correctness dominates:
ACT$\times$aloha's $0.016\to0.824$ swing came entirely from fixing a
normalization bug in our own eval (Table~\ref{tab:act-ablation}), a gap
that dwarfs any cross-architecture difference in the matrix, while
diffusion$\times$pusht measures $\hat{p}{=}0.816$ against the
$0.654$ Hub-card reference~\cite{chi2023diffusion} ($+16.2$\,pp), an
over-count consistent with our last-step success criterion being more
permissive than the card's stricter sticky-\texttt{any(is\_success)}
rule (\texttt{docs/SUCCESS\_CRITERION\_AUDIT.md}). Third, SmolVLA$\times$\texttt{libero\_10}'s 0.252 vs.\
the paper's 0.71 is a scope mismatch
(Section~\ref{sec:results-audit}), not a replication failure. The
contribution we stand behind without qualification is the protocol
itself: a single-PR path for new policies, a shared seed contract, and
CI-vs-MDE bookkeeping that flags sub-threshold deltas as inconclusive.

\subsection{Limitations}

\begin{itemize}\setlength{\itemsep}{0pt}\setlength{\parskip}{0pt}
    \item \textbf{Simulation only.} All numbers are sim; the sim-to-real
    gap is unmeasured. lerobot-bench is a screening tool, not a
    substitute for hardware evaluation.
    \item \textbf{Single-hardware wall-clock.} Per-step latency is from
    one RTX~4060 laptop; absolute ms/step differs across GPUs, but the
    single-GPU \emph{ranking} the auto-downscope rule consumes is portable.
    \item \textbf{$\boldsymbol{\pi_0}$ family deferred.} $\pi_0$,
    $\pi_0$-FAST, $\pi_{0.5}$ (PaliGemma-3B) need ${\approx}30$\,GB host
    RAM during \texttt{from\_pretrained} cold load ($<4$\,GB peak VRAM),
    exceeding the 32\,GB v1 laptop's headroom; v1 compares four learned
    policies, not the targeted seven (v1.1 uses quantized/streaming load).
    \item \textbf{Sparse matrix and episode budget.} Cross-policy
    comparisons are valid only within shared envs (one block: the four
    LIBERO suites). $N{=}250$ is a GPU-budget choice ($N{=}500$ would
    tighten CIs $\sim1.4\times$); headline claims cite only deltas above
    the per-cell MDE.
    \item \textbf{Harness bugs and protocol-vs-paper.} A normalization bug
    in our own eval (self-caught, fixed in PR~\#51;
    Section~\ref{sec:results-audit}) underscores that a harness can be the
    dominant error source; the SmolVLA scope and LIBERO cap mismatches make
    those cells protocol-vs-paper, not replication failures. Cross-policy
    paired tests defer to v1.1 when XVLA returns.
\end{itemize}

\paragraph{Future work.}
Four paths: (1)~close the remaining audit caveats via PR~\#90's
\texttt{--canonical} criterion (cap=600 reruns of the three provisional
LIBERO suites);
(2)~onboard the $\pi_0$ family via quantized/streaming load (four
learned policies $\to$ seven); (3)~add the remaining LIBERO task
variants (10 per suite) and new Hub policies via the single-PR path;
(4)~re-run on hardware for the sim-to-real gap (xvla fixes filed at
\href{https://github.com/huggingface/lerobot/issues/3401}{\texttt{huggingface/lerobot\#3401}}).

% ---------------------------------------------------------------- %
% References                                                       %
% ---------------------------------------------------------------- %

{\small
\bibliographystyle{plain}
\bibliography{references}
}

\end{document}
